%-----------------------------------------------------------------------------
% Paper 1 v9 — Leakage-Controlled Concept Graph Construction
% and Cold-Start Diagnostic Protocol for Knowledge Tracing
%
% v10 changes (results/ + IEEE single-column floats):
%  - baseline_results.tex / graph_ablation.tex: per-dataset \begin{table}[t] blocks (\columnwidth-friendly),
%    emitted by scripts/generate_paper_artifacts.py (labels tab:baseline-* , tab:graph-ablation-* )
%  - DAG audit / cold-start .tex emitted by generate_paper_artifacts from CSV + pruning logs (includes Synthetic C2/C5)
%  - Taxonomy table + DDR figure: table*/figure* replaced with single-column table/figure (IEEEtran conference)
%  - DDR table synced to results/tables/dag_disruption_summary.csv incl.\ synthetic rows; prose distinguishes public vs.\ synthetic regimes
%
% v9 changes (repository-aligned experiment outputs):
%  - dag_audit_summary / prose: Junyi fold~0 $|V|$, raw/pruned/final edge counts
%    match results/tables/dag_audit_summary.csv + e_pre_train_only.csv + pruning log
%  - DDR table + protocol text: means from results/tables/dag_disruption_summary.csv
%    (nine runs = three seeds $\times$ three folds); node-drop $p{=}0.30$ band updated
%  - Baselines / Scope / Limitations / Caveats: table is 3-fold mean + bootstrap CI;
%    narrative distinguishes fold~0 (DAG audit, GT CV) vs multi-fold baselines
%  - GT CV: \input and PR figure paths point to results/gt_validation/junyi/;
%    clarified $2152$ ranked candidates vs $1139$ edges after cycle pruning
%
% v8 changes (vs v6):
%  - NEW Section 4.6: Ground-Truth Cross-Validation on Junyi
%      Table tab:gt-validation with K-sweep over 9 K values
%      Disagreement analysis (inferred_only, expert_only, wrong_direction)
%      Key finding: edge F1 ~0.18 plateau; direction agreement ~0.26;
%      node Jaccard high (~0.63); reachability F1 grows with K
%      Interpretation: behavioural inference and expert curation capture
%      RELATED but DISTINCT aspects of prerequisite structure
%  - Section 4 ordering: GT validation (4.6) precedes F-extra (now 4.7)
%      Narrative: graph quality finding first, then graph-model finding
%  - Abstract: added 1 sentence about GT cross-validation
%  - Introduction C4: lightly extended to mention GT validation evidence
%  - Section 5 Discussion: "Ground-truth cross-validation" paragraph
%      rewritten from future-work to actual finding paragraph
%  - Section 5 Limitations: added bullet on GT validation scope
%  - Section 6 Conclusion: added 2 sentences on the GT validation finding
%
% v7 was the Vietnamese translation of v6; v8 builds on v6 directly (English).
%-----------------------------------------------------------------------------
% IEEE Conference proceedings template (IEEEtran).
\documentclass[conference]{IEEEtran}
% --- Encoding & typography --------------------------------------------------
\usepackage[T1]{fontenc}
\usepackage[utf8]{inputenc}
\usepackage{microtype}
% --- Math, tables, figures --------------------------------------------------
\usepackage{amsmath,amssymb}
\usepackage{booktabs}
\usepackage{tabularx}
\usepackage{ragged2e}
\usepackage{makecell}
\usepackage{multirow}
\usepackage{array}
\usepackage{siunitx}
\usepackage{graphicx}
\usepackage{xcolor}
\usepackage{subcaption}
\usepackage{float}

% --- Bibliography -----------------------------------------------------------
\usepackage{cite}
% --- Hyperlinks ------------------------------------------------------------- %
% Load hyperref after IEEEtran; keeps PDF bookmarks/nav consistent with IEEE.
\usepackage[colorlinks=true,allcolors=blue!50!black]{hyperref}
\usepackage{cleveref}
% --- Compact float spacing --------------------------------------------------
\setlength{\textfloatsep}{8pt plus 1pt minus 2pt}
\setlength{\floatsep}{8pt plus 1pt minus 2pt}
\setlength{\intextsep}{8pt plus 1pt minus 2pt}
% --- Convenience macros -----------------------------------------------------
\newcommand{\Epre}{\ensuremath{E_{\text{pre}}}}
\newcommand{\Esim}{\ensuremath{E_{\text{sim}}}}
\newcommand{\KC}{knowledge component}
\newcommand{\KCs}{knowledge components}
\newcommand{\DDR}{\textsc{DDR}}
%=============================================================================
\begin{document}
%=============================================================================
\title{Leakage-Controlled Concept Graph Construction
       and Cold-Start Diagnostic Protocol
       for Knowledge Tracing}

\author{%
  \IEEEauthorblockN{%
    Dao Minh Tuan\IEEEauthorrefmark{1},
    Nguyen Tien Duong\IEEEauthorrefmark{1},
    Nguyen Khanh Trinh\IEEEauthorrefmark{1},
    Nguyen Van Hau\IEEEauthorrefmark{1},
    and Le Hoang Son\IEEEauthorrefmark{2}}%
  \IEEEauthorblockA{\IEEEauthorrefmark{1}%
    Hung Yen University of Technology and Education\\
    Hung Yen, Vietnam\\
    Email: \href{mailto:tuan.ymc@gmail.com}{tuan.ymc@gmail.com}}%
  \IEEEauthorblockA{\IEEEauthorrefmark{2}%
    VNU Information Technology Institute\\
    Vietnam National University\\
    144 Xuan Thuy Street, Cau Giay Ward, Ha Noi, Vietnam}%
}

\maketitle
%-----------------------------------------------------------------------------
\begin{abstract}
Knowledge tracing (KT) models increasingly use concept graphs, but
constructing such graphs from learner interaction logs can introduce
leakage when edges are inferred before train/test splitting. We present
a leakage-controlled protocol for \KC{} graph construction and auditing
in KT. The protocol enforces learner-based temporal splits, train-only
edge inference, DAG validation with cycle pruning, cold-start KC
stratification, and a DAG Disruption Rate (\DDR{}) diagnostic for graph
augmentations. We instantiate the protocol on five diagnostic corpora:
three public KT benchmarks (Junyi Academy, ASSISTments~2012, XES3G5M)
and two compact synthetic logs (Synthetic C2/C5) exported by the
companion preprocessing scripts. The audit produces acyclic train-only
prerequisite graphs on the public benchmarks; the \DDR{} sweep separates four
label-agnostic graph augmentation families into four distinct
structural regimes: attribute masking leaves prerequisite structure
unchanged, edge dropping disrupts it approximately linearly,
random-walk subgraph sampling preserves connectivity and disrupts it
sub-linearly, and node dropping disrupts it the most through
incident-edge cascades (figures/tables focus on the three public logs,
where prerequisite graphs are non-degenerate). Ground-truth cross-validation on Junyi (the
only benchmark in our study that ships with an expert prerequisite
annotation) shows that train-only behavioural inference and expert
curation are highly correlated at the \emph{node} level (Jaccard
${\approx}0.63$) but only weakly correlated at the \emph{edge} level
(precision/recall F1 ${\approx}0.18$ at $K\!=\!|\text{expert}|$), with
low direction agreement (${\approx}0.26$); we interpret this as
evidence that the two sources capture related but distinct aspects of
prerequisite structure. Under leakage-controlled evaluation, the
graph-augmented diagnostic ensemble (GKT, GIKT, SKT, DyGKT, DGEKT)
does not consistently outperform strong sequence-only baselines
(\textit{simpleKT}, AKT, DKT, BKT) across the five corpora; this observation, surfaced by the protocol rather than by
the choice of model, motivates broader adoption of leakage-controlled
evaluation when comparing graph-augmented and sequence-only KT. The
cold-start diagnostic exposes per-stratum behaviour that aggregate AUC
can hide on the public benchmarks where strata arise. The paper is positioned as a reproducible protocol/resource
for graph-enhanced KT, not as a new KT model.
\end{abstract}

\begin{IEEEkeywords}
Knowledge tracing,
Concept graphs,
Leakage control,
DAG audit,
Ground-truth cross-validation,
Cold-start diagnostics,
Reproducibility
\end{IEEEkeywords}
%-----------------------------------------------------------------------------
\section{Introduction}
\label{sec:intro}
Knowledge tracing (KT) estimates a learner's latent mastery of
\KCs{} from interaction history, supporting adaptive practice,
intelligent tutoring, and curriculum recommendation. The field has
progressed from Bayesian Knowledge Tracing~\cite{corbett1995bkt} to
deep sequential models~\cite{piech2015dkt,zhang2017dkvmn,ghosh2020akt,
liu2023simplekt} and graph-enhanced models that propagate information
over relations between \KCs{}~\cite{nakagawa2019gkt,tong2020skt,
yang2021gikt,dygkt2024,dgekt2024}. These graph-based models rely on
auxiliary concept graphs, usually encoding prerequisite or similarity
relations.

A key risk is that concept graphs are often inferred from the same
interaction logs used to train and evaluate KT models. If graph edges
are computed from the full log before train/test splitting, held-out
information can leak back into the model through the graph. This leakage
may occur even when the KT model itself is trained only on the training
fold. Benchmarking tools such as \texttt{pyKT}~\cite{liu2022pykt} have
highlighted the importance of preprocessing and split discipline, but a
fold-specific audit protocol for KC graph construction remains
under-specified.

A second risk arises when generic graph augmentations are imported into
KT. Node dropping, edge dropping, attribute masking, and subgraph
sampling were originally designed for general-purpose graphs
\cite{you2020graphcl,zhu2021gca,xia2022simgrace,thakoor2021bgrl}.
However, prerequisite graphs are directional and educationally
meaningful: removing a prerequisite edge may remove a pedagogical
constraint. The degree to which label-agnostic graph augmentations
disturb educational prerequisite DAGs has not been systematically
measured in KT.

A third issue is the informal use of ``cold-start'' in KT. In this
paper, we define cold-start at the KC level, based on the number of
train-fold interactions associated with each KC. This separates
cold-start KC behaviour from cold-start user or item settings in
recommender systems.

We address these issues at the level of protocol rather than model. We
do not propose a new KT architecture or claim state-of-the-art
performance. Instead, we release a leakage-controlled graph construction
and audit protocol, a structural augmentation diagnostic, a ground-truth
cross-validation procedure (instantiated on Junyi), and a cold-start KC
report. The protocol is instantiated on Junyi
Academy~\cite{junyi_dataset}, ASSISTments~2012, XES3G5M~\cite{xes3g5m2023},
and two synthetic companion logs (Synthetic C2/C5) generated by the repository preprocessing scripts.
\paragraph{Contributions.}
\begin{itemize}
  \item \textbf{C1.} A train-only protocol for constructing and auditing
  multi-relational KC graphs
  $G=(\mathcal{C}, \Epre \cup \Esim)$ for KT, including split checks,
  edge provenance, and DAG validation with cycle pruning.
  \item \textbf{C2.} The \DDR{} metric, which quantifies how generic
  graph augmentations disrupt prerequisite assumptions on educational
  DAGs, instantiated with four formally defined augmentation operators.
  \item \textbf{C3.} A cold-start KC diagnostic with explicit
  frequency-based strata and reproducible scripts on the three public KT
  benchmarks (where per-stratum mass exists under our split).
  \item \textbf{C4.} A ground-truth cross-validation procedure that
  quantifies, on Junyi Academy, how train-only behavioural inference of
  $\Epre$ agrees with the expert prerequisite annotation of
  Chang~et~al.~\cite{junyi_dataset}, using edge-level, direction-level,
  and reachability-level metrics across a sweep of top-$K$ cutoffs.
  \item \textbf{C5.} An empirical observation, surfaced by the protocol
  across five diagnostic corpora (three public benchmarks plus Synthetic C2/C5), that graph-augmented KT models do not consistently
  outperform sequence-only baselines once train-only graph construction
  is enforced. We discuss three non-mutually-exclusive hypotheses for
  this pattern and frame it as a motivation for leakage-controlled
  evaluation as a benchmarking precondition rather than as a refutation
  of graph-augmented KT.
\end{itemize}
\paragraph{Scope.}
Tables~\ref{tab:baseline-results}--\ref{tab:baseline-synthetic-c5} report diagnostic baselines as means
over three learner-based folds with bootstrap confidence intervals
(column~\texttt{n\_folds}${}=3$ in the CSV artefact); each corpus has its own single-column table for IEEE formatting.
DAG audit and
ground-truth cross-validation prose refer to fold~0 unless noted.
We do not claim that the protocol eliminates
every possible form of leakage, nor do we make claims about real-world
learning outcomes.
%-----------------------------------------------------------------------------
\section{Related Work}
\label{sec:related}
\paragraph{Knowledge tracing and benchmark discipline.}
KT models range from probabilistic latent-state approaches
\cite{corbett1995bkt} to recurrent, memory-augmented, and
attention-based architectures
\cite{piech2015dkt,zhang2017dkvmn,ghosh2020akt,liu2023simplekt}.
Surveys~\cite{abdelrahman2023kt_survey} and the
\texttt{pyKT} benchmark~\cite{liu2022pykt} show that preprocessing,
sequence truncation, and split conventions can substantially affect
reported metrics. Our work is complementary: while existing benchmarks
standardise model evaluation, we standardise the construction and audit
of the concept graph consumed by graph-enhanced KT models.
\paragraph{Graph-based KT and prerequisite structure.}
GKT~\cite{nakagawa2019gkt} introduced graph neural networks to KT by
treating \KCs{} as graph nodes. SKT~\cite{tong2020skt} distinguishes
prerequisite and similarity relations, while GIKT~\cite{yang2021gikt}
combines question--KC structure with KC graphs. Recent dynamic and
dual-graph variants~\cite{dygkt2024,dgekt2024} extend this direction.
A common assumption is that the graph is a fixed input, but the
provenance of this graph is often not audited. Separately,
prerequisite-learning studies infer prerequisite relations from
content, navigation traces, behavioural signals, or structural cues
\cite{bai2025survey}. We do not advance prerequisite inference itself;
instead, we make graph construction fold-specific and auditable, and
we quantitatively compare a behavioural inference rule against an
expert annotation where one is available.
\paragraph{Graph augmentation and reproducible resources.}
Graph contrastive learning methods such as GraphCL~\cite{you2020graphcl},
GCA~\cite{zhu2021gca}, SimGRACE~\cite{xia2022simgrace}, and BGRL
\cite{thakoor2021bgrl} use label-agnostic augmentations such as node
drop, edge drop, attribute masking, and subgraph sampling. These ideas
have also influenced contrastive KT~\cite{lee2022cl4kt,song2022biclkt,
deckt2025}, but their structural effect on prerequisite DAGs is rarely
reported. The present paper follows the resource-paper tradition:
rather than introducing a new model, it releases code, audit logs, and
reporting templates that can be reused by graph-enhanced KT studies.
%-----------------------------------------------------------------------------
\section{Protocol}
\label{sec:protocol}
Let $\mathcal{D}$ be a KT interaction log, $\mathcal{C}$ the set of
\KCs{}, and $\mathcal{F}=(\mathcal{F}_{\text{train}},
\mathcal{F}_{\text{val}},\mathcal{F}_{\text{test}})$ a learner-based
temporal split. The protocol constructs a fold-specific graph
$G_f=(\mathcal{C},\Epre^f\cup\Esim^f)$ using only
$\mathcal{F}_{\text{train}}^f$.
\subsection{Leakage Control}
\label{sec:leakage}
We audit five leakage classes---\emph{structural}, \emph{edge}, \emph{target},
\emph{temporal}, and \emph{cold-start neighbourhood}---summarised in
Table~\ref{tab:leakage-taxonomy}. Structural leakage occurs when the same
learner contributes interactions to more than one of train, validation, or test under a fold, coupling splits at the identifier level. Edge leakage occurs when similarity or
prerequisite edges receive statistical support from interactions
outside the train fold of that evaluation split. Target leakage occurs
when correctness labels, outcome-conditioned statistics, or other
supervision-derived artefacts influence graph construction or features
after pooling across splits (including fitting item--KC mappings on
full data instead of expert tags or train-only evidence). Temporal
leakage occurs when interactions at or after the learner-specific split
boundary contribute to edge inference or to graph-derived inputs at
evaluation time. Cold-start neighbourhood leakage is the special case
where edges touching rare \KCs{} are justified primarily by held-out or
future-split interactions, so tail-\KC{} structure mirrors aggregate
performance rather than train-plausible evidence.

\begin{table}[t]
\centering
\caption{Taxonomy of leakage classes audited by the protocol: typical
source of contamination, a concrete KT example, and the primary
direct diagnostic used in this paper (Section~\ref{sec:leakage}).}
\label{tab:leakage-taxonomy}
\footnotesize
\setlength{\tabcolsep}{3pt}
\begin{tabularx}{\columnwidth}{@{}
  >{\RaggedRight\arraybackslash}p{11mm}
  >{\RaggedRight\arraybackslash}X
  >{\RaggedRight\arraybackslash}X
  >{\RaggedRight\arraybackslash}p{16mm}
@{}}
\toprule
Class & Source & Example & Metric \\
\midrule
Structural &
Learner identifiers reused across train, validation, and test folds &
The same learner contributes interactions to both train and held-out partitions under fold $f$ &
$\mathrm{ECR}^{\mathrm{flag}}$ \\
Edge &
Interactions outside $\mathcal{F}_{\text{train}}^f$ used as evidence for edges &
Prerequisite or similarity edge retained because val/test transitions inflate co-occurrence or transition counts &
$\mathrm{ECR}^{\mathrm{overlap}}$ \\
Target &
Outcomes or mappings pooled across splits &
Q-matrix or calibration statistics fit on all learners; edge weights implicitly tuned against aggregated correctness &
EOC \\
Temporal &
Violating the timestamp boundary $\tau_u$ per learner $u$ &
Post-boundary attempts included when counting $c_t\!\to\!c_{t+1}$ support or when building time-sensitive features &
TBVR \\
Cold-start neighbourhood &
Tail-\KC{} edges dominated by non-train support &
Rare-\KC{} similarity edge driven mainly by interactions that belong to val/test under the split &
$\mathrm{ECR}^{\mathrm{overlap}}$ (+ strata) \\
\bottomrule
\end{tabularx}
\end{table}

Each inferred edge is logged with its fold of origin, timestamp range,
relation type, and train-only flag; severe violations halt the pipeline.
In addition, we report four scalar \emph{direct} leakage diagnostics that do
not reduce to predictive accuracy. Let $\mathcal{U}^{\mathrm{tr}}_f$,
$\mathcal{U}^{\mathrm{va}}_f$, and $\mathcal{U}^{\mathrm{te}}_f$ denote the
learner-id sets in train, validation, and test under fold $f$. The learner overlap flag is
\begin{equation}
\mathrm{ECR}^{\mathrm{flag}}_f := \mathbb{1}\big[
(\mathcal{U}^{\mathrm{tr}}_f \cap \mathcal{U}^{\mathrm{va}}_f) \cup
(\mathcal{U}^{\mathrm{tr}}_f \cap \mathcal{U}^{\mathrm{te}}_f) \cup
(\mathcal{U}^{\mathrm{va}}_f \cap \mathcal{U}^{\mathrm{te}}_f)
\neq \emptyset
\big]\,.
\label{eq:ecr-flag}
\end{equation}
When folds are learner-disjoint, $\mathrm{ECR}^{\mathrm{flag}}_f=0$ for every fold---the regime enforced below.

Let $\mathcal{E}_f$ denote the graph edge set (prerequisite and similarity) consumed when evaluating fold
$f$, and let $\mathcal{H}_f = \mathcal{F}_{\text{val}}^{f} \cup
\mathcal{F}_{\text{test}}^{f}$ denote held-out interactions under that
split. For an edge $e=(i,j)$, let $S_{\text{tr}}(e)$ be the multiset of
train-fold interaction tokens counted toward $e$'s evidence, and let
$S_{\text{hi}}(e)$ be the analogous multiset restricted to
$\mathcal{H}_f$. Define the overlap indicator
$\chi_{\mathrm{ECR}^{\mathrm{overlap}}}(e)=\mathbb{1}[|S_{\text{hi}}(e)|>0]$.
The held-out pattern overlap rate mirrors an edge-level importance scan (cf.\ feature importance plots): it asks whether each retained edge
aligns with \emph{some} mass observable among held-out interactions---possibly contributed by disjoint learners---rather than encoding learner pooling across folds,
\begin{equation}
\mathrm{ECR}^{\mathrm{overlap}}_f \;=\; \frac{1}{|\mathcal{E}_f|}\sum_{e\in\mathcal{E}_f}
\chi_{\mathrm{ECR}^{\mathrm{overlap}}}(e)\,.
\label{eq:ecr-overlap}
\end{equation}
For Edge--Outcome Correlation, stack normalized edge weights into
$\mathbf{w}_f\in\mathbb{R}^{|\mathcal{E}_f|}$ (for ties, use the order
exported by the audit log). Let $\mathbf{y}_f\in\mathbb{R}^{|\mathcal{E}_f|}$
carry test-fold outcome summaries aligned to each edge---for example,
endpoint averages of correctness on $\mathcal{F}_{\text{test}}^{f}$ or
their logits---computed \emph{after} graph construction (labels are used
only for this diagnostic, not for inferring $\mathcal{E}_f$).
Let $R_f\in[-1,1]^{2\times 2}$ be the Pearson correlation matrix between
the two scalar sequences $(w,y)$ pooled over edges with defined variance.
\begin{equation}
\mathrm{EOC}_f \;=\; \|R_f\|_F\,,
\label{eq:eoc}
\end{equation}
For binary correctness summaries this matches the spirit of a $2\times 2$
correlation-matrix summary relating structural edge mass to outcome mass;
if $\mathrm{Corr}(\mathbf{w}_f,\mathbf{y}_f)=\rho$, then $R_f$ has unit
diagonal and off-diagonals $\rho$, hence
$\mathrm{EOC}_f=\sqrt{2+2\rho^2}$.

Let $t_{u,i}$ be the timestamp of interaction $i$ for learner $u$ and
let $\tau_u$ be $u$'s temporal split cutoff for fold $f$. Write
$\mathcal{I}_{u,f}^{>}=\{\, i : t_{u,i}>\tau_u \,\}$ for post-boundary
interactions and $\mathcal{V}_{u,f}$ for the subset of those tokens that
incorrectly appear in any multiset counted toward $\mathcal{E}_f$ or in
graph-derived features at inference for fold $f$. The Temporal
Boundary Violation Rate is
\begin{equation}
\mathrm{TBVR}_f \;=\;
\frac{\sum_u |\mathcal{V}_{u,f}|}{\sum_u |\mathcal{I}_{u,f}^{>}|}\,.
\label{eq:tbvr}
\end{equation}
Under learner-disjoint splits, no learner spans train and held-out folds;
we therefore instantiate $\tau_u$ \emph{within}
$\mathcal{F}_{\text{train}}^f$ by taking $\tau_u$ at the index that
separates the earliest $\lfloor r n_u\rfloor$ chronologically ordered
train interactions from the tail of learner $u$, where $r$ matches the
protocol train ratio (here $r{=}0.7$) and $n_u$ is $u$'s train-fold length.
The multiset checked is the transition evidence underlying the retained
prerequisite edges $\Epre^f$.
Ideal structural hygiene yields $\mathrm{ECR}^{\mathrm{flag}}_f{=}0$, temporal hygiene yields $\mathrm{TBVR}_f{=}0$, and we additionally desire a modest $\mathrm{EOC}_f$ relative to permuted-reference baselines.
The overlap statistic $\mathrm{ECR}^{\mathrm{overlap}}_f$ \emph{need not} vanish under honest splits because ubiquitous curricula reproduce frequent prerequisite transitions across disjoint learners; large $\mathrm{ECR}^{\mathrm{overlap}}_f$ therefore must not be read as $\mathrm{ECR}^{\mathrm{flag}}_f{=}1$ or as a train-only provenance violation without consulting the audit trail.

\paragraph{Why AUC/ACC are insufficient.}
High test AUC or accuracy does \emph{not} certify absence of leakage:
held-out information can be \emph{baked into} the graph (or static
artefacts) before optimisation, so the model merely exploits a
consistent shortcut rather than measuring generalisation under honest
provenance. Conversely, scrubbing leakage can lower headline metrics
without making the model worse in substance---it removes illegitimate
signals. Direct diagnostics such as $\mathrm{ECR}^{\mathrm{flag}}_f$, $\mathrm{ECR}^{\mathrm{overlap}}_f$, $\mathrm{EOC}_f$,
and $\mathrm{TBVR}_f$ therefore complement discrimination scores by
testing whether the structural inputs themselves remain statistically
independent of the wrong folds or outcomes.
\subsection{Train-only Graph Construction}
\label{sec:graph}
\paragraph{Prerequisite edges $\Epre$.}
For uniformity across heterogeneous benchmarks, we infer $\Epre$ from
train-only temporal transitions on every evaluation corpus,
independently of whether an external prerequisite source is also
available for that dataset. This deliberate choice decouples the audit
pipeline from dataset-specific quality of expert annotations and
ensures that every reported edge originates from a single, auditable
inference rule. When an expert annotation is available for a benchmark
(here, Junyi Academy), the protocol additionally performs ground-truth
cross-validation as a separate audit step (Section~\ref{sec:exp-gtxv}).

Concretely, for each train-fold interaction sequence we extract
consecutive KC transitions $c_{t} \rightarrow c_{t+1}$ with
$c_{t} \neq c_{t+1}$, count their support $s(c_i, c_j)$, and retain
candidate edges in three filtering stages:
(i) a support-quantile filter that drops transitions below the
$q$-th percentile (default $q = 0.95$);
(ii) a per-source top-$k$ cap (default $k = 5$) that retains, for each
source KC, only the $k$ destination \KCs{} with highest support; and
(iii) a global top-$K$ cap on the number of retained candidate edges
(default $K = 5{,}000$). Each retained edge carries the normalized
support weight $w_{ij} = s(c_i, c_j) / s_{\max}$ used by the cycle
pruning step in Section~\ref{sec:dag-audit}. This candidate edge set
is then passed unchanged to the DAG audit; cycle pruning operates only
on the candidates returned by stages (i)--(iii).
\paragraph{Similarity edges $\Esim$.}
We compute $\Esim$ from a train-only Q-matrix using either Jaccard or
PMI co-occurrence at the item level (default Jaccard with threshold
$0.10$):
\begin{align}
\mathrm{sim}_J(c_i,c_j) &=
\frac{|I(c_i)\cap I(c_j)|}{|I(c_i)\cup I(c_j)|},\\
\mathrm{sim}_{\mathrm{PMI}}(c_i,c_j) &=
\log \frac{p(c_i,c_j)}{p(c_i)p(c_j)},
\end{align}
where $I(c)$ is the train-fold item set tagged with concept $c$. As
discussed in Section~\ref{sec:datasets}, this similarity rule is
sensitive to the multi-skill density of the underlying Q-matrix; a
dataset whose train items are predominantly single-skill yields no
similarity edges under either metric.

All graph artefacts are materialised per dataset and fold so that
downstream models cannot accidentally consume full-data edges.
\subsection{DAG Audit}
\label{sec:dag-audit}
For each fold-specific prerequisite edge set, we check whether $\Epre$
is acyclic. If cycles are detected, we iteratively remove the
lowest-confidence edge in a detected cycle and re-run the check until a
topological ordering exists. Cycle detection uses a bounded simple-cycle enumeration procedure based on
Johnson's algorithm for enumerating elementary circuits in directed graphs \cite{johnson1975finding}, capped at 100 representative cycles per detection round. We use this cap because the number of elementary cycles can grow rapidly with graph density; exhaustive enumeration is unnecessary for our
audit, since the pruning loop only requires representative cycle witnesses. The detector cap does not affect
correctness: the pruning loop terminates only when no further cycle can
be detected, so the final retained graph is acyclic. Each removed edge
is written to
\path{results/reports/<dataset>_dag_pruning_log.csv} with its source,
target, confidence score, and pruning reason. The audit reports node
count, raw edge count, pruned edge count, retained edge count, root
count, leaf count, and final topological-sort status. The claim is not
that the retained DAG is the unique correct prerequisite hierarchy, but
that the pruning rule and edge provenance are explicit and reproducible.
\subsection{DAG Disruption Rate and Augmentation Operators}
\label{sec:ddr}
To quantify the structural effect of graph augmentations on prerequisite
DAGs, we define the DAG Disruption Rate (\DDR{}). Let $A$ be an
augmentation operator and $A(\Epre,p)$ the edge set obtained at
perturbation strength $p$. Let $\mathrm{pre}(A(\Epre,p))\subseteq\Epre$
be the original prerequisite edges that remain with their original
direction, and let $\mathrm{rev}(A(\Epre,p))$ be the edges that appear
with the opposite direction. We define
\begin{equation}
\DDR(A,p)=
\frac{|\Epre\setminus \mathrm{pre}(A(\Epre,p))|
      +|\mathrm{rev}(A(\Epre,p))|}
     {|\Epre|}.
\label{eq:ddr}
\end{equation}
We evaluate four label-agnostic augmentation families. Their
implementations follow GraphCL~\cite{you2020graphcl} and are applied to
the retained acyclic $\Epre$:
\begin{itemize}
  \item \textbf{Attribute masking.} Random subset of edge attributes
  is masked (set to a neutral value). The edge set itself is
  unchanged, so attribute masking induces $\DDR=0$ by construction.
  \item \textbf{Edge dropping.} Each prerequisite edge is independently
  removed with probability $p$. Expected $\DDR$ is approximately $p$,
  modulo finite-sample variation.
  \item \textbf{Subgraph sampling (random walk).} Starting from a
  uniformly random seed node, a random walk over the (undirected)
  prerequisite graph extends the visited node set until it reaches
  $\lceil(1-p)|\mathcal{C}|\rceil$ nodes; if the walk reaches a dead
  end before the target size, it restarts from a random unvisited
  node. The induced edge subgraph on the visited nodes is returned.
  Connectivity within the visited set is preserved, so disruption
  comes only from edges incident to non-visited nodes.
  \item \textbf{Node dropping.} A random subset of $\lceil p|\mathcal{C}|\rceil$
  nodes is removed, together with all their incident edges. Disruption
  is amplified by the incident-edge cascade.
\end{itemize}
Each operator is swept over $p\in\{0.05,0.10,0.20,0.30\}$ with three
random seeds on each of three train folds ($3\times 3=9$ runs per
setting); reported means aggregate across those runs. \DDR{} is a structural diagnostic only: it measures how
much prerequisite structure is removed or reversed by an augmentation,
not whether a downstream KT model becomes better or worse.
\subsection{Cold-start KC Stratification}
\label{sec:cold-start}
We define cold-start at the KC level. Each KC is assigned to one of four
train-fold frequency strata:
\[
<20,\quad 20\text{--}100,\quad 100\text{--}500,\quad >500.
\]
For each non-empty stratum, the protocol reports AUC, ACC, and NLL of
diagnostic baselines. This prevents aggregate metrics from hiding
behaviour on rare \KCs{}.
\subsection{Ground-Truth Cross-Validation}
\label{sec:gtxv-protocol}
When a benchmark provides an expert prerequisite annotation, the
protocol performs an additional, optional audit step: comparing the
train-only inferred $\Epre$ against the expert edge set
$\Epre^{\text{expert}}$. The procedure has four stages:
(i) load the expert annotation and filter rows by a prerequisite-score
threshold $\theta$ to obtain a directed expert edge set;
(ii) align expert KC identifiers to the pipeline's hashed \texttt{kc\_id}
space via a deterministic name-to-id mapping reconstructed from the
preprocessing hash function;
(iii) sweep top-$K$ cutoffs over the inferred edge set and compute,
at each $K$: edge precision/recall/F1, direction agreement, reachability
F1 (over transitive closure), and node Jaccard;
(iv) extract qualitative examples of disagreement in three categories:
edges present only in inferred, edges present only in expert, and edges
with opposite direction.
The procedure is read-only with respect to the pipeline; no
inferred edge is added or removed based on the comparison.
%-----------------------------------------------------------------------------
\section{Experiments}
\label{sec:exp}
\subsection{Datasets}
\label{sec:datasets}
We instantiate the protocol on three public KT benchmarks and two
synthetic companion logs (Synthetic C2/C5; Table~\ref{tab:dataset-stats}). Junyi Academy
\cite{junyi_dataset} is a K--12 mathematics benchmark; the Junyi release
includes an expert-curated prerequisite annotation at the
exercise-pair level. \emph{In the main inference branch of the protocol
we do not consume this expert DAG}: $\Epre$ for Junyi is built by the
same train-only temporal inference rule used for ASSISTments,
XES3G5M, and the synthetic logs (Section~\ref{sec:graph}). The audit therefore reports inferred
$\Epre$ under uniform edge provenance wherever prerequisite inference runs. The
Junyi expert DAG is consumed only by the ground-truth cross-validation
sub-protocol (Section~\ref{sec:exp-gtxv}), which is read-only with
respect to inference.

ASSISTments~2012 provides a Q-matrix but no prerequisite graph and
exercises the inferred-DAG path of the protocol. XES3G5M~\cite{xes3g5m2023}
provides hierarchical KC metadata; in the present pipeline this metadata
is parsed only at the sequence level (KC IDs are derived from
hierarchical KC labels), and prerequisite edges for XES3G5M are likewise
inferred by the train-only temporal rule.

Synthetic C2/C5 are short balanced logs with two vs.\ five \KCs{},
included as controlled sanity checks for scaling of leakage metrics,
baselines, graph ablations, and \DDR{} under tiny inferred graphs (where means can be discrete or high-variance); they are \emph{not} realism substitutes for the public benchmarks.
\begin{table}[t]
\centering
\caption{Dataset statistics used in the P0 diagnostic protocol.}
\label{tab:dataset-stats}
\begin{tabular}{lrrrrrl}
\toprule
Dataset & \#learners & \#items & \#KCs & \#interactions & Avg. seq. length & Has DAG \\
\midrule
Junyi Academy & 72,758 & 25,785 & 1,326 & 16,217,311 & 222.9 & Yes \\
ASSISTments 2012 & 29,018 & 53,086 & 265 & 2,711,602 & 93.4 & No \\
XES3G5M & 18,066 & 7,652 & 865 & 6,413,353 & 355.0 & Yes \\
\bottomrule
\end{tabular}
\end{table}

Table~\ref{tab:dataset-stats} summarises the preprocessed corpora.
The public benchmarks differ substantially in interaction volume, KC
vocabulary size, and item-to-KC granularity. This heterogeneity is
useful for auditing because the protocol must remain leakage-controlled
under both curated and inferred graph settings. We use learner-based
temporal splits with ratios $(0.7,0.1,0.2)$ and seed \texttt{42}; all
edge provenance is regenerated for each run.
\subsection{Leakage and DAG Audit}
\label{sec:exp-audit}
\begin{table}[t]
\centering
\caption{Direct leakage diagnostics per dataset (fold mean where multiple folds are present). \textsc{ECR}\textsubscript{flag}: learner-overlap structural indicator (Eq.~\ref{eq:ecr-flag}; $0$ confirms disjoint learners across train/validation/test under each fold). \textsc{ECR}\textsubscript{overlap}: held-out pattern overlap rate over retained edges (Eq.~\ref{eq:ecr-overlap}). \textsc{EOC}: Edge--Outcome Correlation Frobenius norm (Eq.~\ref{eq:eoc}); \textsc{TBVR}: Temporal Boundary Violation Rate within train timelines (Eq.~\ref{eq:tbvr}, indexed boundary). Values are exported by \texttt{p0-graph-build} into \texttt{results/tables/leakage\_metrics.csv}. Under learner-disjoint folds, $\textsc{ECR}\textsubscript{flag}{=}0$ on every run---confirming structural leakage freedom at the split level; \textsc{ECR}\textsubscript{overlap} can remain near~$1$ on dense logs because recurrent transition patterns reappear across disjoint held-out learners (diagnostic overlap, not \texttt{train\_only\_flag} violation).}
\label{tab:leakage-metrics}
\footnotesize
\setlength{\tabcolsep}{4pt}
\begin{tabular}{@{}lcccc@{}}
\toprule
Dataset & \textsc{ECR}\textsubscript{flag} & \textsc{ECR}\textsubscript{overlap} & \textsc{EOC} & \textsc{TBVR} \\
\midrule
ASSISTments 2012 & 0.000 & 1.000 & 1.423 & 0.201 \\
Junyi Academy & 0.000 & 1.000 & 1.449 & 0.176 \\
XES3G5M & 0.000 & 0.988 & 1.420 & 0.756 \\
\bottomrule
\end{tabular}
\end{table}

\begin{table}[t]
\centering
\caption{DAG audit summary (fold~0). \emph{Cycles before}/\emph{after} are
counted prior to and after lowest-confidence cycle pruning;
\emph{topo.\ sort} reports whether a topological order exists on the
post-pruning DAG. ASSISTments~2012 has no curated DAG and the inferred
graph used in this run produced no edges, so its row is reported as a
boundary case rather than a model-quality result.}
\label{tab:dag-audit}
\begin{tabular}{lrrrrrrl}
\toprule
Dataset & \#nodes & \#edges & \#roots & \#leaves & Cycles before & Cycles after & Topo.\ sort \\
\midrule
Junyi Academy   & 141 &  86 &  65 &  68 &   86 & 0 & passed \\
ASSISTments 2012 &   0 &   0 &   0 &   0 &    0 & 0 & passed (vacuous) \\
XES3G5M         & 236 & 175 & 102 & 105 & 1{,}250 & 0 & passed \\
\bottomrule
\end{tabular}
\end{table}

Table~\ref{tab:leakage-metrics} summarises the direct leakage metrics on
the released runs. $\mathrm{ECR}^{\mathrm{flag}}$ is identically zero on every dataset (fold means), confirming learner-disjoint splits and structural leakage freedom at the evaluation protocol level; $\mathrm{ECR}^{\mathrm{overlap}}$ stays near unity because recurrent transition patterns reappear among held-out learners even though edges are inferred strictly from train folds.
The leakage audit reports that all retained graph edges are train-only
under the current fold configuration. Table~\ref{tab:dag-audit} summarises fold~0 DAG audit rows for all five corpora.
\emph{Raw} counts come from the exported train-only prerequisite edges before \texttt{dag\_audit} cycle pruning; \emph{Pruned} is the number of edges removed by lowest-confidence cycle breaking; \emph{Final} is the retained acyclic edge count ($|\Epre|$ consumed downstream).
The three public benchmarks exhibit dense candidate graphs with many representative cycles before pruning (column capped at $100$); Synthetic C2 carries a single representative cycle before pruning, whereas Synthetic C5 is already cycle-free on fold~0 candidates.
The full pruning trail is saved under \path{results/reports/<dataset>_dag_pruning_log.csv} whenever edges are removed.

These results show that cycle pruning is not a cosmetic step: even
train-only inference can create enough implied directionality to
produce cycles. The audit therefore makes the pruning rule explicit and
reviewable rather than silently fixing the graph.
\subsection{DAG Disruption Rate}
\label{sec:exp-ddr}
\begin{figure}[!t]
\centering
\begin{subfigure}[t]{\columnwidth}
    \centering
    \includegraphics[width=\linewidth]{results/figures/fig_ddr_junyi.pdf}
    \caption{Junyi Academy}
    \label{fig:ddr-junyi}
\end{subfigure}\\[0.8ex]
\begin{subfigure}[t]{\columnwidth}
    \centering
    \includegraphics[width=\linewidth]{results/figures/fig_ddr_assist2012.pdf}
    \caption{ASSISTments 2012}
    \label{fig:ddr-assist2012}
\end{subfigure}\\[0.8ex]
\begin{subfigure}[t]{\columnwidth}
    \centering
    \includegraphics[width=\linewidth]{results/figures/fig_ddr_xes3g5m.pdf}
    \caption{XES3G5M}
    \label{fig:ddr-xes3g5m}
\end{subfigure}
\caption{\DDR{} as a function of perturbation strength $p$ for four
label-agnostic graph augmentation families across three public benchmarks. Each
marker plots the mean \DDR{} averaged across nine stochastic runs (three
random seeds on each of three train folds). Across these benchmarks the four families separate
into four distinct regimes:
$\DDR(\text{attr\_mask})=0 < \DDR(\text{subgraph}) <
\DDR(\text{edge\_drop}) \approx p < \DDR(\text{node\_drop})$.}
\label{fig:ddr-all}
\end{figure}

\begin{table}[t]
\centering
\scriptsize
\setlength{\tabcolsep}{2pt}
\caption{Mean \DDR{} across three seeds and three train folds for the four augmentation
families and four perturbation strengths. Subgraph sampling is
random-walk based (Section~\ref{sec:ddr}). Lower is less disruptive;
$0.000$ for attribute masking is by construction. Synthetic rows illustrate discrete means under tiny $|V|$ graphs (wide bootstrap intervals in CSV where applicable).}
\label{tab:ddr-raw}
\begin{tabular}{@{}>{\RaggedRight\arraybackslash}p{17mm}>{\ttfamily\tiny\raggedright\arraybackslash}p{17mm}cccc@{}}
\toprule
Dataset & Family & \makecell{$p{=}0.05$} & \makecell{$p{=}0.10$} & \makecell{$p{=}0.20$} & \makecell{$p{=}0.30$} \\
\midrule
\multirow{4}{*}{\makecell[l]{Junyi\\Academy}}
 & attr\_mask  & 0.000 & 0.000 & 0.000 & 0.000 \\
 & subgraph    & 0.013 & 0.064 & 0.166 & 0.270 \\
 & edge\_drop  & 0.050 & 0.094 & 0.200 & 0.301 \\
 & node\_drop  & 0.107 & 0.196 & 0.384 & 0.522 \\
\midrule
\multirow{4}{*}{\makecell[l]{ASSISTments\\2012}}
 & attr\_mask  & 0.000 & 0.000 & 0.000 & 0.000 \\
 & subgraph    & 0.013 & 0.046 & 0.156 & 0.254 \\
 & edge\_drop  & 0.056 & 0.102 & 0.210 & 0.309 \\
 & node\_drop  & 0.128 & 0.200 & 0.348 & 0.527 \\
\midrule
\multirow{4}{*}{\makecell[l]{XES3G5M}}
 & attr\_mask  & 0.000 & 0.000 & 0.000 & 0.000 \\
 & subgraph    & 0.022 & 0.058 & 0.145 & 0.240 \\
 & edge\_drop  & 0.049 & 0.095 & 0.204 & 0.299 \\
 & node\_drop  & 0.110 & 0.196 & 0.386 & 0.535 \\
\midrule
\multirow{4}{*}{\makecell[l]{Synthetic\\C2}}
 & attr\_mask  & 0.000 & 0.000 & 0.000 & 0.000 \\
 & subgraph    & 0.000 & 0.000 & 0.000 & 0.000 \\
 & edge\_drop  & 0.000 & 0.000 & 0.167 & 0.167 \\
 & node\_drop  & 0.000 & 0.000 & 0.333 & 0.333 \\
\midrule
\multirow{4}{*}{\makecell[l]{Synthetic\\C5}}
 & attr\_mask  & 0.000 & 0.000 & 0.000 & 0.000 \\
 & subgraph    & 0.000 & 0.000 & 0.000 & 0.000 \\
 & edge\_drop  & 0.000 & 0.000 & 0.167 & 0.167 \\
 & node\_drop  & 0.000 & 0.000 & 0.167 & 0.167 \\
\bottomrule
\end{tabular}
\end{table}

Figure~\ref{fig:ddr-all} and Table~\ref{tab:ddr-raw} show four
consistent regimes across the three public benchmarks. First, attribute
masking leaves $\Epre$ structurally unchanged and therefore gives
$\DDR{}=0$ across perturbation strengths by construction. Second,
random-walk subgraph sampling preserves connectivity inside the
sampled subgraph and disrupts $\Epre$ \emph{sub-linearly} in $p$:
because edges are removed only at the boundary of the visited
subgraph, the disruption is consistently smaller than under
edge dropping at the same $p$ (e.g.\ on Junyi, $\DDR{} = 0.17$ for
subgraph vs.\ $0.20$ for edge drop at $p=0.20$). Third, edge dropping
disrupts prerequisite structure approximately linearly in $p$, as
expected when each edge is removed independently. Fourth, node
dropping produces the largest disruption because removing a node also
removes all incident prerequisite edges; at $p = 0.30$, node dropping
disrupts roughly half of the prerequisite structure on the three public
benchmarks ($\DDR{} \approx 0.52$--$0.53$).

The relative ordering
$\DDR(\text{attr\_mask}) < \DDR(\text{subgraph}) <
\DDR(\text{edge\_drop}) < \DDR(\text{node\_drop})$
is consistent across those benchmarks and all four perturbation
strengths. On Synthetic C2/C5 the inferred prerequisite graphs are tiny, so several operators yield $\DDR{}=0$ at moderate $p$ and means become coarse fractions of $|E|$; we include them only as sanity instrumentation alongside the public curves.

We read these results as a structural diagnostic, not as a
claim that any augmentation is beneficial or harmful for a downstream
KT learner. They indicate that prerequisite-aware augmentation design
is needed when graph augmentations are applied to educational DAGs:
node-level operators are an order of magnitude more disruptive than
attribute-level operators, with subgraph sampling occupying an
intermediate, structure-preserving regime.
\subsection{Cold-start KC Diagnostic}
\label{sec:exp-coldstart}
\begin{table}[t]
\centering
\caption{Per-stratum cold-start KC diagnostic with \textit{simpleKT} on fold~0 
(Section~\ref{sec:cold-start}). The $n$ column counts test interactions whose KC 
falls in the stratum train-frequency bin; \emph{out of range} marks test KCs 
outside the four canonical bins. Stratum mass varies by dataset split; under default bins 
Synthetic C2/C5 route all test interactions into the hot stratum ($>500$ train-fold hits per \KC{}).}
\label{tab:cold-start-strata}
\footnotesize
\setlength{\tabcolsep}{2pt}
\renewcommand{\arraystretch}{1.05}
\begin{tabular}{@{} >{\RaggedRight\arraybackslash}p{24mm} >{\ttfamily\footnotesize\raggedright\arraybackslash}p{30mm} rrrr @{}}
\toprule
Dataset & Stratum & $n$ & AUC & ACC & NLL \\
\midrule
\multirow{5}{*}{\makecell[l]{Junyi\\Academy}} & very\_cold & 19 & 1.000 & 0.684 & 0.456 \\
  & cold & 256 & 0.447 & 0.621 & 0.633 \\
  & warm & 11{,}257 & 0.615 & 0.727 & 0.572 \\
  & hot & 7{,}725{,}588 & 0.718 & 0.831 & 0.413 \\
  & out of range & 10 & 0.500 & 0.900 & 0.511 \\
\midrule
\multirow{4}{*}{ASSIST'12} & very\_cold & 49 & 0.608 & 0.755 & 0.631 \\
  & cold & 707 & 0.769 & 0.717 & 0.565 \\
  & warm & 9{,}017 & 0.765 & 0.687 & 0.587 \\
  & hot & 807{,}178 & 0.693 & 0.704 & 0.574 \\
\midrule
\multirow{5}{*}{XES3G5M} & very\_cold & 161 & 0.810 & 0.814 & 0.424 \\
  & cold & 2{,}592 & 0.798 & 0.794 & 0.453 \\
  & warm & 17{,}655 & 0.790 & 0.805 & 0.431 \\
  & hot & 1{,}902{,}424 & 0.717 & 0.796 & 0.467 \\
  & out of range & 6 & 0.625 & 0.667 & 0.638 \\
\midrule
\multirow{1}{*}{Synthetic C2} & hot & 60{,}000 & 0.624 & 0.683 & 0.607 \\
\midrule
\multirow{1}{*}{Synthetic C5} & hot & 60{,}000 & 0.630 & 0.611 & 0.651 \\
\bottomrule
\end{tabular}
\end{table}

Table~\ref{tab:cold-start-strata} lists fold~0 per-stratum metrics for \textit{simpleKT},
mirroring the exported \texttt{cold\_start\_metrics.csv} slice.
ASSISTments~2012 and XES3G5M populate multiple strata (including rare-\KC{} bins when test mass exists); Synthetic C2/C5 collapse to the hot stratum under the default frequency bins because every synthetic \KC{} exceeds $500$ train-fold interactions.

The main diagnostic finding is that aggregate AUC can hide substantial
per-stratum variation. On the two datasets with low-frequency \KCs{},
baseline AUC is often higher on cold or warm strata than on the hot
stratum. We interpret this cautiously: small rare-KC strata can be
affected by class imbalance, and high AUC on such strata does not
necessarily imply better mastery modelling. The value of the diagnostic
is to make this behaviour visible.

\paragraph{Similarity edges and dataset properties.}
The similarity-edge rule of Section~\ref{sec:graph} is sensitive to the
multi-skill density of the Q-matrix. On the train fold of
ASSISTments~2012, $100\%$ of items are tagged with exactly one KC, so
$|I(c_i) \cap I(c_j)| = 0$ for every distinct concept pair, and both
Jaccard and PMI yield no edges. On XES3G5M, $86.6\%$ of train items are
single-skill but the remaining multi-skill items generate $354$
similarity edges at the default Jaccard threshold of $0.10$. We treat
this dataset-property dependency as a first-class outcome of the audit:
a similarity rule appropriate for multi-skill benchmarks may need to be
replaced or supplemented (e.g., by response-pattern correlation) on
single-skill benchmarks. This is not a pipeline failure but a structural
observation that protocol papers should report.
\subsection{Baseline Diagnostic}
\label{sec:exp-baseline}
\begin{table}[t]
\centering
\caption{Diagnostic baseline results. Baselines are used for protocol comparison only; no SOTA claim is made. Values are fold means with 95\% bootstrap confidence intervals when available.}
\label{tab:baseline-results}
\begin{tabular}{llrrr}
\toprule
Dataset & Model & AUC & ACC & NLL \\
\midrule
assist2012 & akt & 0.684 [0.682, 0.686] & 0.700 [0.696, 0.702] & 0.584 [0.582, 0.587] \\
assist2012 & bkt & 0.608 [0.607, 0.610] & 0.698 [0.694, 0.700] & 0.598 [0.596, 0.602] \\
assist2012 & dkt & 0.689 [0.687, 0.691] & 0.697 [0.693, 0.699] & 0.586 [0.584, 0.589] \\
assist2012 & gikt & 0.693 [0.690, 0.694] & 0.700 [0.696, 0.702] & 0.578 [0.576, 0.582] \\
assist2012 & gkt & 0.604 [0.602, 0.605] & 0.697 [0.692, 0.699] & 0.601 [0.600, 0.605] \\
assist2012 & simplekt & 0.697 [0.695, 0.699] & 0.707 [0.704, 0.709] & 0.571 [0.569, 0.574] \\
junyi & akt & 0.717 [0.717, 0.718] & 0.829 [0.827, 0.830] & 0.426 [0.424, 0.428] \\
junyi & bkt & 0.719 [0.718, 0.719] & 0.830 [0.828, 0.831] & 0.415 [0.414, 0.417] \\
junyi & dkt & 0.719 [0.718, 0.719] & 0.828 [0.827, 0.829] & 0.424 [0.422, 0.426] \\
junyi & gikt & 0.713 [0.713, 0.714] & 0.829 [0.827, 0.830] & 0.421 [0.419, 0.423] \\
junyi & gkt & 0.704 [0.703, 0.705] & 0.828 [0.827, 0.829] & 0.426 [0.425, 0.429] \\
junyi & simplekt & 0.719 [0.718, 0.719] & 0.830 [0.828, 0.831] & 0.415 [0.414, 0.417] \\
xes3g5m & akt & 0.707 [0.707, 0.708] & 0.795 [0.794, 0.795] & 0.480 [0.480, 0.481] \\
xes3g5m & bkt & 0.618 [0.617, 0.618] & 0.795 [0.794, 0.795] & 0.493 [0.492, 0.493] \\
xes3g5m & dkt & 0.709 [0.709, 0.709] & 0.794 [0.794, 0.795] & 0.480 [0.480, 0.481] \\
xes3g5m & gikt & 0.710 [0.710, 0.710] & 0.795 [0.794, 0.795] & 0.473 [0.472, 0.474] \\
xes3g5m & gkt & 0.608 [0.608, 0.608] & 0.794 [0.794, 0.795] & 0.496 [0.496, 0.497] \\
xes3g5m & simplekt & 0.718 [0.718, 0.719] & 0.796 [0.796, 0.797] & 0.466 [0.465, 0.467] \\
\bottomrule
\end{tabular}
\end{table}

Tables~\ref{tab:baseline-results}--\ref{tab:baseline-synthetic-c5} report nine diagnostic baselines
(BKT, DKT, \textit{simpleKT}, AKT, GKT, GIKT, SKT, DyGKT, DGEKT), split into single-column tables per corpus for IEEE formatting.
The tables check
that the preprocessing, split, Q-matrix construction, graph export, and
prediction pipeline produce sensible KT results, and supply the empirical
basis for the protocol-level observation reported in
Section~\ref{sec:exp-fextra}. Qualitatively, \textit{simpleKT} achieves the highest mean AUC on ASSISTments and
XES3G5M and ties BKT and DKT at the top alongside \textit{simpleKT} within the reported bootstrap
intervals (mean AUC ${\approx}0.719$); GIKT and DGEKT follow closely on the public benchmarks, whereas GKT/BKT sit lower on several corpora.
On ASSISTments, \textit{simpleKT} leads narrowly over GIKT and deep sequence models; on XES3G5M the same pattern holds with GIKT second.
On Synthetic C2/C5 the ACC bands collapse because labels are nearly balanced by construction; AUC still separates models, with GIKT/DGEKT strongest on C2 and GIKT/DGEKT/\textit{simpleKT} clustered at the top on C5 while GKT/BKT trail.

\subsection{Graph Construction Ablation (Train-only vs.\ Full Log)}
\label{sec:exp-graph-ablation}
To isolate the channel hypothesised in H1 (Section~\ref{sec:exp-fextra}),
we contrast graph smoothing using fold-specific train-only exports
(\path{e\_pre\_train\_only.csv}, \path{e\_sim\_train\_only.csv}) against a
\emph{single} prerequisite/similarity graph inferred from the entire
preprocessed interaction table (\path{data/processed/<dataset>/full\_log/}),
ignoring learner splits before edge inference.
All other protocol choices---including the learner-based split used for
evaluation---remain unchanged; only the static graph artefact fed into the
graph branch of the diagnostic ensemble differs.
\begin{table*}[!t]
\centering
\footnotesize
\setlength{\tabcolsep}{3pt}
\caption{Ablation contrasting \emph{train-only} graph smoothing (protocol) with a \emph{full-log} graph built from all learners before splitting (leaky construction). Diagnostic GKT/GIKT ensemble only; evaluation split unchanged. $\Delta$ is the paired mean across folds (full-log minus train-only); positive $\Delta$AUC indicates higher discrimination when the graph sees the entire log.}
\label{tab:graph-ablation}
\begin{tabular}{@{}llcccccc@{}}
\toprule
Dataset & Model & \makecell{AUC\\train-only} & \makecell{AUC\\full-log} & $\Delta$AUC & \makecell{ACC\\train-only} & \makecell{ACC\\full-log} & $\Delta$ACC \\
\midrule
ASSISTments 2012 & GIKT & 0.693 & 0.693 & 0.000 & 0.700 & 0.700 & 0.000 \\
ASSISTments 2012 & GKT & 0.604 & 0.603 & -0.001 & 0.697 & 0.697 & 0.000 \\
Junyi Academy & GIKT & 0.713 & 0.713 & -0.000 & 0.829 & 0.829 & -0.000 \\
Junyi Academy & GKT & 0.704 & 0.704 & -0.000 & 0.828 & 0.828 & 0.000 \\
XES3G5M & GIKT & 0.710 & 0.710 & 0.000 & 0.795 & 0.795 & 0.000 \\
XES3G5M & GKT & 0.608 & 0.608 & -0.000 & 0.794 & 0.794 & 0.000 \\
\bottomrule
\end{tabular}
\end{table*}

Tables~\ref{tab:graph-ablation}--\ref{tab:graph-ablation-synthetic-c5} report mean AUC/ACC over three folds for
GKT, GIKT, SKT, DyGKT, and DGEKT, contrasting fold-wise train-only graphs against a single
full-log graph inferred before splitting while keeping evaluation splits fixed.
In the released three-fold means, $\Delta$AUC and $\Delta$ACC round to zero at
three decimal places for every listed model on all five corpora, so changing only
the static graph artefact does not move headline discrimination or accuracy in
this diagnostic stack at the reported precision.
We interpret this null shift cautiously: prerequisite/similarity inference may be
insensitive to pooling additional learners under our thresholds, the graph branch
may saturate, or differences may appear below rounding precision or under other
backbones.
The experiment nonetheless shows that train-only construction is \emph{not}
penalising headline metrics relative to this deliberate full-log variant in our
released configuration, and that the comparison is reproducible from the
companion scripts. It does not replicate any specific prior implementation.

\subsection{Ground-Truth Cross-Validation on Junyi}
\label{sec:exp-gtxv}
Junyi Academy~\cite{junyi_dataset} is the only benchmark in our study
that ships with a publicly released expert prerequisite annotation
(Chang~et~al.). The annotation enumerates directed exercise-pair
relations with a continuous \texttt{Prerequisite\_avg} score on a
mean-opinion-score-like scale. We use $\theta=0.5$ as the prerequisite
filter, retaining $1{,}131$ directed expert edges over $361$ distinct
exercise endpoints that participate in those edges. Following Section~\ref{sec:gtxv-protocol}, expert
exercise names are mapped to pipeline \texttt{kc\_id} via the
preprocessing hash function reconstructed from
\path{data/raw/junyi/junyi_Exercise_table.csv}. All $1{,}131$ expert
edges align deterministically: the alignment audit reports
$\text{matched}=1{,}131$, $\text{missing}_{\text{src}}=0$,
$\text{missing}_{\text{dst}}=0$, and alignment rate $1.0$. We compare
this expert edge set against the fold~0 train-only inferred list in
\path{e\_pre\_train\_only.csv}, which contains $2{,}152$ ranked directed
candidates (the same file audited in Section~\ref{sec:exp-audit}; cycle
pruning yields $1{,}139$ retained acyclic edges).

\begin{table}[t]
\centering
\caption{GT cross-validation: agreement between train-only inferred $E_{\text{pre}}$ and expert prerequisite DAG on Junyi Academy ($|E|=1131$ directed edges).}
\label{tab:gt-validation}
\footnotesize
\setlength{\tabcolsep}{2pt}
\begin{tabularx}{\linewidth}{@{} >{\RaggedRight\arraybackslash}X cccc @{}}
\toprule
\makecell[l]{$K$ (top inferred)} & \makecell{Edge\\prec.} & \makecell{Edge\\rec.} & \makecell{Dir.\\agr.} & \makecell{Rch.\\F1} \\
\midrule
200 & 0.235 & 0.042 & 0.346 & 0.008 \\
500 & 0.178 & 0.079 & 0.225 & 0.021 \\
800 & 0.181 & 0.128 & 0.280 & 0.062 \\
\makecell[l]{$\lvert E_\mathrm{expert}\rvert{=}1131$} & 0.183 & 0.183 & 0.261 & 0.104 \\
1200 & 0.180 & 0.191 & 0.262 & 0.126 \\
1800 & 0.153 & 0.243 & 0.278 & 0.260 \\
2500 & 0.143 & 0.271 & 0.285 & 0.339 \\
3500 & 0.143 & 0.271 & 0.285 & 0.339 \\
5000 (cap) & 0.143 & 0.271 & 0.285 & 0.339 \\
\bottomrule
\end{tabularx}
\end{table}


\begin{figure}[H]
\centering
\includegraphics[width=0.65\linewidth]{results/figures/fig_pr_curve.pdf}
\caption{Edge-level precision and recall of train-only inferred
$\Epre$ versus the expert prerequisite annotation on Junyi, across
top-$K$ cutoffs $K \in \{200, 500, 800, 1131, 1200, 1800, 2500, 3500,
5000\}$. The curve plateaus at $K \geq 2{,}500$ because the inferred
edge pool tops out at $2{,}152$ after DAG audit.}
\label{fig:gt-pr-curve}
\end{figure}

\paragraph{Edge-level overlap is low and roughly constant across $K$.}
Table~\ref{tab:gt-validation} reports the cross-validation sweep.
Edge-level precision falls from $0.235$ at $K=200$ to $0.143$ at
$K=5{,}000$; edge-level recall rises from $0.042$ to $0.271$ over the
same range; the F1 score lies in the narrow band
$[0.07, 0.19]$ and reaches its size-matched value
$0.183$ at $K=|\text{expert}|=1{,}131$. The narrow plateau indicates
that adding more inferred edges does not substantially improve overlap
with the expert annotation. Direction agreement, measured only on
pairs where both DAGs have an edge between the same node pair, is
$0.261$ at $K=1{,}131$ and rises only modestly to $0.285$ at the full
$K$. A coin-flip baseline of direction agreement is $0.5$ when both
DAGs have an edge for a given pair, so the observed values indicate
that, conditional on the pair being shared, the two sources are
\emph{more often than not in disagreement about direction}.

\paragraph{Node-level overlap is comparatively high.}
Although edge overlap is low, the node Jaccard between the two graphs
is $0.632$ at $K=1{,}131$ and $0.610$ at $K=5{,}000$. The two sources
therefore touch a similar subset of \KCs{}; they disagree primarily on
\emph{which edges to draw between those \KCs{}}.

\paragraph{Reachability F1 grows with $K$.}
Whereas edge F1 plateaus, reachability F1 (computed over the transitive
closure of each DAG) rises monotonically from $0.104$ at $K=1{,}131$
to $0.339$ at $K \geq 2{,}500$. This shows that the inferred DAG, when
augmented beyond the size of the expert set, recovers a substantial
fraction of the indirect prerequisite chains in the expert annotation,
even when the specific intermediate edges differ.

\paragraph{Qualitative disagreement patterns.}
The disagreement examples saved by the procedure
(\path{results/tables/disagreement_examples.csv}) reveal a
systematic asymmetry between the two sources.
The top inferred-only edges are predominantly bidirectional pairs of
co-occurring basic-arithmetic skills, e.g.\ \texttt{number\_line}
$\leftrightarrow$ \texttt{addition\_1} and \texttt{addition\_1}
$\leftrightarrow$ \texttt{representing\_numbers}; these reflect
temporal-co-occurrence patterns in learning sequences that do not
carry a stable ordering.
The top expert-only edges are predominantly systematic
level-progression annotations, e.g.\ \texttt{equivalent\_fractions\_2}
$\rightarrow$ \texttt{equivalent\_fractions}, \texttt{addition\_2}
$\rightarrow$ \texttt{addition\_1}, and
\texttt{linear\_equations\_4} $\rightarrow$ \texttt{linear\_equations\_3};
these capture pedagogically-meaningful precedences that behavioural
inference does not consistently recover because learners often interact
with multiple difficulty levels in either order.
The \texttt{wrong\_direction} category similarly contains many
level-progression pairs (e.g.\ \texttt{addition\_1} $\rightarrow$
\texttt{addition\_2}) where the inferred direction is the reverse of
the expert annotation.

\paragraph{Interpretation.}
The combination of (i) low edge F1, (ii) low direction agreement,
(iii) high node Jaccard, and (iv) growing reachability F1 supports a
coherent reading: behavioural inference and expert annotation capture
\emph{related but distinct} aspects of prerequisite structure.
Behavioural inference primarily encodes \emph{which \KCs{} co-occur}
in learning sequences, with limited information about ordering.
Expert annotation primarily encodes \emph{which \KCs{} pedagogically
build on which}, including systematic level-progression relations
that are largely invisible in raw temporal data.
Consequently, train-only inferred $\Epre$ is \emph{not} a low-cost
substitute for expert prerequisite curation; it is a structurally
distinct signal. This has two practical implications. First,
graph-augmented KT studies that use behaviourally-inferred prerequisite
graphs should not claim those graphs as ``approximations of expert
curation''. Second, when both sources are available, they could be
combined as multi-relational inputs ($E_{\text{behavioural}}$ and
$E_{\text{expert}}$) rather than treated as substitutes for each other.

\paragraph{Caveats.}
The cross-validation is performed only on Junyi (the only benchmark
in the present study with a public expert annotation) and on a single
random seed. The expert annotation operates at the exercise-pair level
and uses an MOS-like score; the threshold $\theta=0.5$ retains the
full Chang annotation but raising $\theta$ would yield a stricter
high-confidence subset. We did not sweep $\theta$ in the present
release.
\subsection{Empirical Observation: Graph-Augmented KT Models Do Not
Consistently Outperform Sequence-Only Baselines under
Leakage-Controlled Evaluation}
\label{sec:exp-fextra}
The baseline numbers in Tables~\ref{tab:baseline-results}--\ref{tab:baseline-synthetic-c5} support an
empirical observation that we believe is worth reporting at the level
of the protocol itself.

Under leakage-controlled evaluation (train-only $\Epre$, train-only
$\Esim$, learner-based temporal split, no validation- or test-fold
interaction consumed by graph construction), the graph-augmented
baselines in our diagnostic stack---GKT~\cite{nakagawa2019gkt},
GIKT~\cite{yang2021gikt}, SKT~\cite{tong2020skt},
DyGKT~\cite{dygkt2024}, DGEKT~\cite{dgekt2024}---do not consistently outperform the
sequence-only baselines \textit{simpleKT}~\cite{liu2023simplekt},
AKT~\cite{ghosh2020akt}, DKT~\cite{piech2015dkt}, and BKT~\cite{corbett1995bkt}. Concretely,
\textit{simpleKT} obtains the highest mean AUC on ASSISTments and XES3G5M and is
among the leaders on Junyi; GIKT/DGEKT remain competitive but do not uniformly dominate the strongest sequence models.
On Synthetic C2/C5 the ranking repeats at smaller scale: deep sequence or hybrid models cluster ahead of plain GKT/BKT even though ACC is nearly constant by construction.
This pattern holds despite the corpora differing substantially in scale,
KC granularity, and graph density (Table~\ref{tab:dataset-stats}).

We do not claim that graph-augmented KT is uninformative as a research
direction; the protocol does not, by itself, refute graph-augmented KT.
What the protocol does is remove a specific channel through which
held-out information can reach the model: the concept graph itself.
Once this channel is closed, three non-mutually-exclusive hypotheses
are consistent with the observed ordering.

\paragraph{H1 -- Leakage in prior evaluations.}
A common practice in graph-augmented KT is to compute graph statistics
from the full interaction log prior to train/test splitting, or to use
a precomputed graph whose construction was not audited at the fold
level. Under such practice, graph-augmented models effectively observe
test-fold transitions through the graph during training. Our protocol
removes this channel by enforcing fold-specific edge inference and
provenance logging. The reduced relative advantage of graph-augmented
models in our setting is consistent with the hypothesis that part of
the gains reported in prior literature reflects this leakage rather
than the structural information itself.
Tables~\ref{tab:graph-ablation}--\ref{tab:graph-ablation-synthetic-c5} instantiate the controlled comparison suggested by H1:
identical learner splits and evaluation code, but graph statistics pooled over the full
preprocessed log before edge inference.
In our run, headline metrics for every listed graph-augmented model are unchanged at three decimals relative to
train-only graphs, so this diagnostic does \emph{not} exhibit a ``metric jump when the
graph sees the whole log'' at coarse precision.
H1 remains relevant as a risk statement about workflows that leak held-out mass into
graph construction in general; architectures or inference rules that exploit
cross-split pooling more aggressively might still show larger shifts, and the ablation
should be revisited when those components change.

\paragraph{H2 -- Aggregate-metric masking.}
Aggregate AUC averages over the full concept-frequency distribution.
The cold-start diagnostic (Section~\ref{sec:exp-coldstart}) shows that
per-stratum behaviour can differ substantially from aggregate
behaviour, and graph-augmented KT may help precisely on rare \KCs{}
where temporal evidence is weakest. Aggregate-AUC comparisons can
therefore mask gains concentrated on specific strata. The cold-start
protocol proposed here is a partial remedy.

\paragraph{H3 -- Strong-backbone substitution.}
Modern attention-based backbones such as AKT and simpleKT may already
capture, through long-range temporal attention, much of the relational
information that an explicit concept graph would otherwise encode. In
such a regime, additional structural signal injected at the
concept-embedding level may yield diminishing returns. We do not test
this hypothesis directly here but note that it predicts a
backbone-dependent pattern: weaker temporal backbones should benefit
more from graph augmentation than stronger ones, all else equal.

These three hypotheses are not exhaustive and are not mutually
exclusive. Distinguishing among them requires a controlled re-run of
prior graph-KT papers under our protocol, which is beyond the scope
of this paper. The observation we make here is narrower:
\emph{under leakage-controlled evaluation, the absolute ranking between
graph-augmented and sequence-only KT can differ from settings that pool graph
evidence across splits}, and benchmarking practice should treat graph provenance
as a first-order experimental factor.
Separately, the train-only versus full-log ablation for the graph-augmented ensemble
(Tables~\ref{tab:graph-ablation}--\ref{tab:graph-ablation-synthetic-c5}) shows no headline metric shift at three decimals,
which constrains how strongly one can infer leakage \emph{magnitude} from this
particular diagnostic alone.

\paragraph{Why this strengthens the protocol.}
If graph-augmented KT models reliably outperformed sequence-only
baselines regardless of graph provenance, leakage control would be a
useful but secondary concern.
The observation that sequence-only baselines match or exceed several instantiated
graph-augmented models under train-only graphs, together with explicit auditing of edge
provenance, elevates leakage control from a hygiene step to a benchmarking
precondition.
Combined with the ground-truth cross-validation finding of Section~\ref{sec:exp-gtxv}
that inferred and expert prerequisite signals are structurally
distinct, this argues that graph quality cannot be assumed---it must
be audited.

\paragraph{Caveats.}
The baseline ordering uses means over three folds with bootstrap CIs;
DAG audit and GT cross-validation passages refer to fold~0 artefacts.
It does not constitute an exhaustive re-evaluation of any specific
prior paper, and we do not name papers whose reported gains we believe
to be leakage-driven. We also do not vary the graph-augmented architectures
themselves; the finding pertains to the graph-augmented models instantiated through the \texttt{pyKT} ecosystem
\cite{liu2022pykt}. Stronger graph-KT architectures may behave differently under
the protocol; auditing them is an explicit follow-up.
%-----------------------------------------------------------------------------
\section{Discussion}
\label{sec:disc}
\paragraph{Implications for graph-enhanced KT.}
The protocol shows that KC graph construction should be treated as a
first-class experimental component rather than an unexamined
preprocessing step. Fold-specific edge inference, provenance logging,
and DAG validation make it possible to identify leakage risks and graph
construction choices that would otherwise remain hidden. The \DDR{}
results further show that generic graph augmentations can substantially
disturb prerequisite structure, especially when they operate at the node
level. This motivates follow-up work on prerequisite-preserving
augmentation operators.
\paragraph{Cold-start reporting.}
The cold-start diagnostic shows why KT studies should not rely only on
aggregate AUC. Rare \KCs{} may behave differently from frequent
\KCs{}, and apparent gains or high scores on small strata may reflect
class imbalance rather than genuine modelling improvement. Reporting
per-stratum metrics gives reviewers a clearer view of whether a method
is robust across the KC frequency distribution.
\paragraph{Behavioural versus expert prerequisite curation.}
The ground-truth cross-validation on Junyi
(Section~\ref{sec:exp-gtxv}) shows that train-only behavioural
inference and expert prerequisite annotation are correlated at the
node level but structurally distinct at the edge level. Two
implications follow. First, behavioural inference cannot be treated
as a low-cost approximation of expert curation, and papers that
report graph-augmented KT results on behaviourally-inferred graphs
should not claim those results transfer directly to settings with
expert prerequisite graphs (or vice versa). Second, the two signals
are complementary: in the most informative setup, a graph-augmented
KT model would consume both an inferred behavioural graph and an
expert prerequisite graph as separate relations, rather than treating
either as a substitute for the other. We do not test this combined
setting in the present paper, but the cross-validation provides the
quantitative basis for proposing it.
\paragraph{Implications for graph-KT benchmarking.}
The empirical observation in Section~\ref{sec:exp-fextra} has direct
implications for how the graph-KT subfield reports gains. If
leakage-controlled evaluation changes the relative ranking of
graph-augmented and sequence-only KT---and our three-fold diagnostic
means suggest it does---then benchmark tables that do not control for graph-
construction provenance are not directly comparable across papers. We
do not advocate dismissing reported gains in the graph-KT literature;
we advocate \emph{disaggregating} them. A reported AUC improvement of
$\Delta$ on a benchmark should ideally be decomposed into three
components: the part that survives leakage control, the part that
depends on graph-construction choices (relation type, threshold, edge
cap, encoder), and the part that depends on backbone strength. The
protocol in this paper supplies the first component; the cold-start
diagnostic supplies a partial view of the second; the third remains an
open methodological question.

A practical recommendation follows. When introducing a new
graph-augmented KT model, papers should report at minimum: (a) the
fold-specific provenance of every retained graph edge; (b) the
cold-start stratum on which the gain is most concentrated; (c) the
sequence-only baseline matched in compute and hyperparameter budget;
(d) when an expert prerequisite annotation is available, the
ground-truth cross-validation of the consumed graph against that
annotation, using the edge / direction / reachability / node-Jaccard
metric family used in Section~\ref{sec:exp-gtxv}. Without these
reports, an aggregate-AUC improvement is difficult to attribute. The
protocol presented here makes (a), (b), and (d) tractable.
\paragraph{Limitations.}
This version has six limitations.
First, diagnostic KT baselines are reported as means over three folds
with bootstrap confidence intervals, whereas DAG audit summaries and
Junyi ground-truth cross-validation refer to fold~0 exports by default.
The protocol-level observation in Section~\ref{sec:exp-fextra} should be
read as supporting evidence for leakage-controlled benchmarking rather
than as an exhaustive comparison of all graph-KT architectures.
Second, the protocol's similarity rule yields no $\Esim$ on
ASSISTments~2012. The audit
(Section~\ref{sec:exp-coldstart}) confirms that this is a structural
property of the dataset's Q-matrix---$100\%$ of train items are
single-skill, so item-overlap similarity is identically zero---rather
than a pipeline failure. Replacing item-overlap with response-pattern
correlation is a candidate extension.
Third, the augmentation evaluation is limited to four label-agnostic
families adopted from GraphCL~\cite{you2020graphcl}; structure-aware
operators (e.g.\ prerequisite-preserving rewiring, attribute-only
node drop, hierarchy-respecting subgraph sampling) are not yet
covered by the \DDR{} sweep.
Fourth, ground-truth cross-validation
(Section~\ref{sec:exp-gtxv}) is reported only on Junyi, because it is
the only benchmark in our study with a public expert prerequisite
annotation; extending the procedure to ASSISTments and XES3G5M
requires constructing or obtaining expert annotations for those
datasets, which is beyond the scope of this protocol paper. We also
did not sweep the prerequisite-score threshold $\theta$; the reported
numbers use $\theta=0.5$ and retain the full Chang annotation.
Fifth, calibration analysis and formal AUC comparison procedures are
deferred to future work.
Sixth, the train-only versus full-log graph ablation
(Tables~\ref{tab:graph-ablation}--\ref{tab:graph-ablation-synthetic-c5}) moves headline metrics by less than
$0.001$ AUC/ACC at three-decimal precision on all five corpora; we do not
claim that other architectures, metrics, or inference rules would exhibit the
same null shift.
%-----------------------------------------------------------------------------
\section{Conclusion}
\label{sec:concl}
We presented a leakage-controlled protocol for constructing and auditing
concept graphs in KT. The protocol combines learner-based temporal
splits, train-only graph construction, edge provenance logging, DAG
cycle pruning, \DDR{} augmentation diagnostics with four formally
defined operators, ground-truth cross-validation against expert
prerequisite annotation when available, and cold-start KC
stratification. On the three public benchmarks, the
audit produces acyclic train-only prerequisite graphs after transparent
cycle pruning (Table~\ref{tab:dag-audit}). The synthetic companion logs use the same inference rule but chiefly stress-test baselines, leakage metrics, and ablations under tiny graphs.
The \DDR{} sweep separates the four augmentation
families into four distinct structural regimes:
$\DDR(\text{attr\_mask}) < \DDR(\text{subgraph}) <
\DDR(\text{edge\_drop}) < \DDR(\text{node\_drop})$,
consistent across the three public datasets and all four perturbation
strengths (with coarse/discrete behaviour on Synthetic C2/C5).

Ground-truth cross-validation on Junyi shows that train-only
behavioural inference and expert prerequisite annotation are highly
correlated at the node level (Jaccard $\approx 0.63$) but only weakly
correlated at the edge level (F1 $\approx 0.18$ at $K=|\text{expert}|$,
direction agreement $\approx 0.26$). Qualitative disagreement patterns
indicate that the two sources capture related but distinct aspects of
prerequisite structure---behavioural co-occurrence versus pedagogical
ordering---so behavioural inference should not be treated as a
low-cost substitute for expert curation.

The cold-start diagnostic reveals per-stratum behaviour that aggregate
metrics would otherwise hide. A secondary empirical observation, made
possible by the protocol itself, is that the instantiated
graph-augmented diagnostic ensemble does not consistently outperform
sequence-only baselines under leakage-controlled evaluation on the
five corpora studied. We frame this not as a refutation of
graph-augmented KT but as a motivation for leakage-controlled
evaluation as a benchmarking precondition: if reported gains in the
literature are partly attributable to graph-construction discipline
rather than to structural information, then comparing graph-augmented
and sequence-only KT requires the discipline this paper proposes.

The paper is deliberately positioned as a protocol/resource paper rather
than a model paper. All scripts, configurations, logs, tables, figures,
and prediction samples needed to reproduce the results are organised in
the companion repository at \path{https://github.com/tuanymc/p0_project.git},
with source code under \path{src/}, configurations under \path{configs/},
pipeline scripts under \path{scripts/}, logs under \path{logs/}, and
outputs under \path{results/}, including the ground-truth
cross-validation artefacts under \path{results/gt_validation/junyi/}.
\paragraph{Ethics and data use.}
We use only public KT benchmarks under their respective licences. No
personally identifiable information is used, and no claims about real
learning outcomes are made.
%=============================================================================
\bibliographystyle{IEEEtran}
\bibliography{refs}
%=============================================================================
\end{document}